\documentclass[11pt]{article}

% ── NeurIPS 2024 style ─────────────────────────────────────────────────────

% ── Core packages ──────────────────────────────────────────────────────────
\usepackage[utf8]{inputenc}
\usepackage[T1]{fontenc}
\usepackage{microtype}
\usepackage{hyperref}
\usepackage{url}
\usepackage{booktabs}
\usepackage{multirow}
\usepackage{multicol}
\usepackage{amsmath}
\usepackage{amssymb}
\usepackage{graphicx}
\usepackage{xcolor}
\usepackage{enumitem}
\usepackage{array}
\usepackage{tabularx}
\usepackage{makecell}
\usepackage{rotating}
\usepackage{colortbl}
\usepackage{caption}
\usepackage{subcaption}
\usepackage{float}
\usepackage{wrapfig}
\usepackage{tcolorbox}
\usepackage{mdframed}
\usepackage{longtable}
\usepackage{geometry}
\usepackage{setspace}
\usepackage[numbers]{natbib}

% ── Colour palette ─────────────────────────────────────────────────────────
\definecolor{headerbg}{RGB}{220,235,248}
\definecolor{altrow}{RGB}{240,245,250}
\definecolor{accentblue}{RGB}{31,78,145}
\definecolor{lightgray}{RGB}{245,245,245}
\definecolor{darkgray}{RGB}{60,60,60}
\definecolor{indicoprange}{RGB}{204,85,0}
\definecolor{successgreen}{RGB}{0,128,0}

% ── tcolorbox styles ──────────────────────────────────────────────────────
\tcbuselibrary{skins,breakable}
\newtcolorbox{abstractbox}{
  enhanced,
  breakable,
  colback=lightgray,
  colframe=accentblue!60,
  arc=2pt,
  boxrule=0.6pt,
  left=6pt, right=6pt, top=4pt, bottom=4pt,
  title={\textbf{Abstract}},
  fonttitle=\small\bfseries,
  coltitle=accentblue,
  attach boxed title to top left={yshift=-2mm, xshift=6pt},
  boxed title style={colback=white, colframe=accentblue!60, arc=2pt, boxrule=0.4pt}
}

\newtcolorbox{researchbox}[1]{
  enhanced,
  breakable,
  colback=indicoprange!5,
  colframe=indicoprange!50,
  arc=2pt,
  boxrule=0.6pt,
  left=6pt, right=6pt, top=4pt, bottom=4pt,
  title={#1},
  fonttitle=\small\bfseries,
  coltitle=indicoprange,
}

% ── Hyperref setup ─────────────────────────────────────────────────────────
\hypersetup{
  colorlinks=true,
  linkcolor=accentblue,
  citecolor=accentblue,
  urlcolor=accentblue,
  pdftitle={Tool Calling in LLMs: A Survey from English-Centric Benchmarks to Indic Frontiers},
  pdfauthor={Anonymous}
}

% ── Column type for tables ────────────────────────────────────────────────
\newcolumntype{L}[1]{>{\raggedright\arraybackslash}p{#1}}
\newcolumntype{C}[1]{>{\centering\arraybackslash}p{#1}}
\newcolumntype{R}[1]{>{\raggedleft\arraybackslash}p{#1}}

% ── Misc macros ───────────────────────────────────────────────────────────
\newcommand{\taubench}{$\tau$-bench}
\newcommand{\tautwobench}{$\tau^2$-bench}
\newcommand{\bfcl}{\textsc{bfcl}}
\newcommand{\yes}{\textcolor{successgreen}{\checkmark}}
\newcommand{\no}{\textcolor{red}{$\times$}}
\newcommand{\partialmark}{\textcolor{orange}{$\sim$}}
\newcommand{\eg}{\textit{e.g.},\xspace}
\newcommand{\ie}{\textit{i.e.},\xspace}
\newcommand{\cf}{\textit{cf.}\xspace}

% ══════════════════════════════════════════════════════════════════════════
\title{%
  Tool Calling in Large Language Models:\\
  A Survey from English-Centric Benchmarks\\
  to Indic Frontiers%
}

\author{%
  Sai Teja M S\\
  AI4Bharat, IIT Madras\\
  \texttt{me21b171@smail.iitm.ac.in}
}

% ══════════════════════════════════════════════════════════════════════════
\begin{document}

\maketitle

% ── ABSTRACT ──────────────────────────────────────────────────────────────
\begin{abstractbox}
Tool calling --- the ability of Large Language Models (LLMs) to invoke external APIs, functions, and
databases in response to natural language queries --- has become the foundational capability for agentic
AI systems. Since Toolformer~\citep{schick2023toolformer} demonstrated self-supervised tool learning
and Gorilla~\citep{patil2023gorilla} showed that fine-tuned models can surpass GPT-4 on API call
accuracy, the field has rapidly expanded to cover parallel function calls, multi-turn agent
interactions, policy-adherent workflows, and real-world deployment benchmarks. This survey provides
the first systematic overview of the tool-calling literature across three dimensions:
(1)~\textbf{foundational approaches} from early prompting to reinforcement learning;
(2)~\textbf{evaluation benchmarks} from stateless single-turn tests to stateful dual-control
environments; and (3)~\textbf{the critical gap in Indic language coverage}, where despite representing
over 1.4 billion speakers across 22 scheduled languages, tool-calling research remains almost entirely
English-centric. We survey 45+ papers spanning 2022--2026, introduce the newly proposed International
Tool Calling (ITC) dataset~\citep{itc2025} into the multilingual picture, and present a concrete
research roadmap for Indic tool calling, including a benchmark specification
(\textsc{IndicFuncBench}) covering 10 Indic languages and four India-specific API domains.
\end{abstractbox}

\vspace{4pt}

% ══════════════════════════════════════════════════════════════════════════
\section{Introduction}
\label{sec:intro}
% ══════════════════════════════════════════════════════════════════════════

\subsection{Motivation and Scope}

Large Language Models are increasingly deployed not as standalone conversational systems but as
\emph{agents} --- systems that perceive their environment through API calls and natural language,
reason about goals, and take actions with real-world consequences. The capability enabling this shift
is \textbf{tool calling} (also called \emph{function calling} or \emph{tool use}): the model
generates structured calls to external functions, databases, web services, or specialised APIs,
receives results, and uses them to compose a response or plan further actions.

The economic stakes are high. Agentic AI is being deployed in customer service (booking,
cancellations, refunds), enterprise workflow automation, government citizen services, and healthcare
triage. Whether a model correctly calls \texttt{cancel\_reservation(id=X)} versus
\texttt{modify\_reservation(id=X)} --- or knows to ask for clarification before calling either ---
has direct financial and safety consequences.

This survey covers the tool-calling literature from 2022 to April 2026, spanning foundational
methods, evaluation frameworks, training paradigms, and the nascent but critical area of multilingual
and Indic-language tool calling.

\subsection{Why an Indic Extension Matters}

India has 1.4 billion people, 22 constitutionally scheduled languages, and hundreds of dialects.
Hindi alone has 600 million speakers; Bengali, Tamil, Telugu, Marathi, Kannada, Gujarati, Odia,
Malayalam, and Punjabi each have tens to hundreds of millions. Yet the dominant benchmarks remain
strikingly English-centric:

\begin{itemize}[leftmargin=1.2em, itemsep=2pt]
  \item \textbf{\bfcl{}} --- the dominant function-calling leaderboard --- is English-only except
        for Java/JavaScript programming language tests~\citep{patil2025bfcl}.
  \item \textbf{\taubench{}} and \textbf{\tautwobench{}} contain English-only user simulators and
        task annotations~\citep{yao2024taubench,barres2025tau2}.
  \item \textbf{MASSIVE-Agents} (ACL 2025) includes Hindi but limits Indic script coverage;
        non-Hindi Indic languages are largely absent~\citep{massiveagents2025}.
  \item Indic LLMs (Sarvam-1, Krutrim, Airavata, OpenHathi) have no dedicated tool-calling
        benchmarks.
  \item Real deployments --- government citizen services (Aadhaar, IRCTC, NPCI), UPI payments,
        e-commerce (Flipkart, Meesho), healthcare (Practo, PharmEasy) --- urgently need agents that
        handle tool calling in Hindi, Tamil, Telugu, Bengali, and other Indic languages.
\end{itemize}

\begin{researchbox}{Positioning Statement}
This survey is the first to systematically connect the global tool-calling literature with the Indic
NLP ecosystem. Prior tool-calling surveys~\citep{wangsurvey2024,springersurvey2025} are
English-centric; prior Indic LLM surveys~\citep{indicanalysis2025} cover no tool calling. We bridge
this gap with a complete survey and a concrete research roadmap.
\end{researchbox}

% ══════════════════════════════════════════════════════════════════════════
\section{A Taxonomy of Tool Calling}
\label{sec:taxonomy}
% ══════════════════════════════════════════════════════════════════════════

We propose a five-dimensional taxonomy (Table~\ref{tab:taxonomy}) that unifies the fragmented
vocabulary used across the literature.

\begin{table}[h]
\centering
\caption{Five-dimensional taxonomy of tool calling in LLMs.}
\label{tab:taxonomy}
\renewcommand{\arraystretch}{1.25}
\begin{tabular}{L{2.8cm} L{9.8cm}}
\toprule
\rowcolor{headerbg}
\textbf{Dimension} & \textbf{Values} \\
\midrule
\textbf{D1: Interaction Mode}
  & Single-turn stateless $\cdot$ Single-turn multi-call $\cdot$
    Multi-turn conversational $\cdot$ Agentic multi-step \\
\rowcolor{altrow}
\textbf{D2: Tool Selection}
  & Given single tool $\cdot$ Select from small pool $\cdot$
    Retrieve from large pool (1000+ APIs) $\cdot$ Dynamic tool creation \\
\textbf{D3: Control Structure}
  & Single-control (agent only) $\cdot$
    Dual-control (agent + user, Dec-POMDP) \\
\rowcolor{altrow}
\textbf{D4: Evaluation}
  & AST matching $\cdot$ Executable evaluation $\cdot$
    State-based evaluation $\cdot$ LLM-as-judge \\
\textbf{D5: Training}
  & Prompting-only (zero/few-shot) $\cdot$ SFT $\cdot$
    Retriever-Aware Training (RAT) $\cdot$ Reinforcement learning \\
\bottomrule
\end{tabular}
\end{table}

\paragraph{D1 -- Interaction Mode.}
The simplest setting is \emph{single-turn stateless}: one natural-language query produces one
structured call. \emph{Single-turn multi-call} settings (parallel function calling) require
simultaneous independent calls. \emph{Multi-turn} settings carry state across turns, while fully
\emph{agentic} settings have the model taking autonomous sequential actions without per-step user
input.

\paragraph{D2 -- Tool Selection Complexity.}
Selection complexity ranges from trivially filling arguments for a single given tool, through
choosing among a handful of candidates with distractors, to retrieving the correct API from pools of
thousands (ToolBench: 16,000+ APIs from RapidAPI Hub~\citep{qin2024toolllm}). The most demanding
setting requires the model to \emph{synthesise} a new tool specification when no existing tool
satisfies the query.

\paragraph{D3 -- Control Structure.}
Standard benchmarks assume \emph{single-control}: only the agent has tools.
\tautwobench{}~\citep{barres2025tau2} introduced the \emph{dual-control} setting formalised as a
Decentralised Partially Observable Markov Decision Process (Dec-POMDP), where both the agent and the
user possess separate tool sets acting on a shared world state. The best model achieves only 49\%
pass@1 in this setting versus 61--80\% in single-control benchmarks.

% ══════════════════════════════════════════════════════════════════════════
\section{Foundational Methods: A Chronological Survey}
\label{sec:methods}
% ══════════════════════════════════════════════════════════════════════════

\subsection{Pre-2023: Early Tool Augmentation}

\paragraph{TALM \citep{parisi2022talm}.}
One of the first systematic studies showing that augmenting LLMs with tool calls significantly
expands capability. Used a self-play loop to generate training data for a fixed set of calculators
and search tools.

\paragraph{LaMDA \citep{thoppilan2022lamda}.}
Google's dialogue model demonstrated early tool use (calculator, search) via fine-tuning on
tool-augmented demonstrations, establishing that structured tool calls can be embedded in
conversational turns.

\paragraph{ReAct \citep{yao2023react}.}
Proposed interleaving \emph{Thought} (reasoning trace) and \emph{Action} (tool call) in LLM output.
The Thought $\rightarrow$ Action $\rightarrow$ Observation chain substantially outperforms
Act-only prompting and remains the de facto prompting paradigm for all subsequent agent benchmarks.

\paragraph{WebShop \citep{yao2022webshop}.}
Simulated an e-commerce environment where agents must browse web pages and execute purchase actions,
establishing the grounded language-agent paradigm.

\subsection{2023: The Grammar of Tool Use}

\paragraph{Toolformer \citep{schick2023toolformer}.}
Introduced self-supervised learning of when and how to call tools. The model learns to insert API
calls as special tokens by bootstrapping: it generates candidate call positions, executes the calls,
and retains those that reduce perplexity on the continuation. While limited to a small fixed set
(calculator, search, Wikipedia, calendar, translation), it demonstrated that tool use can be learned
without human annotation of call positions.

\paragraph{Gorilla \citep{patil2023gorilla}.}
Introduced Retriever-Aware Training (RAT), fine-tuning LLaMA-7B to write accurate API calls from
HuggingFace, TorchHub, and TensorHub documentation. On APIBench (1,645 ML API calls), Gorilla
surpasses GPT-4 on API call accuracy, especially with a document retriever handling API version
changes. Crucially, Gorilla introduced Abstract Syntax Tree (AST) comparison for measuring API
hallucination --- a methodology later adopted by \bfcl{}.

\paragraph{HuggingGPT \citep{shen2023hugginggpt}.}
Used ChatGPT as a controller to plan, dispatch, and integrate calls to specialised HuggingFace
models as tools, establishing the ``LLM as orchestrator of specialist models'' paradigm.

\paragraph{API-Bank \citep{li2023apibank}.}
Introduced a three-level evaluation hierarchy: (1) Call, (2) Retrieve+Call, and (3)
Plan+Retrieve+Call, across 73 frequently-used APIs. GPT-4 excels at planning; GPT-3.5 struggles at
Plan+Retrieve+Call.

\paragraph{ToolLLM / ToolBench \citep{qin2024toolllm}.}
Fine-tuned LLaMA on 16,000+ real-world APIs from RapidAPI Hub using a novel Depth-First Search
Decision Tree (DFSDT) reasoning strategy allowing backtracking across solution paths. With 120,000+
instruction-API pairs, ToolBench catalysed the tool-learning fine-tuning ecosystem.

\paragraph{NexusRaven \citep{srinivasan2023nexusraven}.}
A commercially permissive (Apache 2.0) function-calling model. NexusRaven-13B was the first
open-source model to claim surpassing GPT-4 on zero-shot function calling benchmarks using a
multi-step data generation pipeline.

\paragraph{MetaTool \citep{huang2023metatool}.}
Tested whether LLMs can decide \emph{when} to use tools (not just how), establishing the
relevance/irrelevance detection evaluation axis later adopted by \bfcl{}.

\paragraph{ToolAlpaca \citep{tang2023toolalpaca}.}
Used a multi-agent simulation environment to automatically generate a tool-use corpus with
$\sim$3,000 simulated tools, demonstrating generalised tool learning without needing real API access.

\subsection{2024: Scale, Multi-Turn, and Reliability}

\paragraph{BFCL v1$\rightarrow$v3 \citep{patil2025bfcl}.}
The Berkeley Function Calling Leaderboard evolved from a static AST-evaluated benchmark
(v1: 2,000 examples) to a live community-contribution dataset (v2: 2,251 examples with
contamination detection) to a full multi-turn state-based evaluation framework (v3: 5,551 pairs
across Base, Missing-Params, Missing-Funcs, Long-Context, and Composite categories). No model
exceeds $\sim$75\% overall on v3.

\paragraph{\taubench{} \citep{yao2024taubench}.}
Introduced LM-simulated users, domain policy compliance, and the \texttt{pass}$^k$ reliability
metric (probability all $k$ trials succeed). GPT-4o achieves 61\% pass@1 in $\tau$-retail but
below 25\% pass@8, exposing dramatic consistency failures at scale. Trajectory analysis reveals:
wrong argument (33\%), wrong decision/policy violation (25\%), wrong info (22\%), partial
resolution (19\%).

\paragraph{ToolACE \citep{liu2025toolace}.}
Introduced Tool Self-Evolution Synthesis (TSS), generating diverse APIs from scratch using
speciation-adaptation-evolution, and Self-Guided Dialog Generation (SDG) for complexity
calibration. ToolACE-8B achieves competitive \bfcl{} scores versus GPT-4, setting the state of
the art for 8B open-source models.

\paragraph{ToolSandbox \citep{lu2025toolsandbox}.}
Introduced stateful tool execution with implicit state dependencies between tools, on-policy
conversational evaluation, and a dynamic milestone/minefield evaluation strategy. Reveals a
significant open-source vs.\ proprietary performance gap on stateful tasks.

\paragraph{xLAM \citep{liu2024xlam}.}
A family of Large Action Models fine-tuned on diverse tool-use datasets across 21 domains,
establishing a scalable fine-tuning pipeline for open-source tool-using agents.

\paragraph{StableToolBench \citep{guo2024stabletoolbench}.}
Addresses API instability in ToolBench with a virtual API server implementing caching and
LLM-powered simulation for reproducible evaluation.

\paragraph{HammerBench \citep{wang2024hammerbench}.}
Fine-grained function-calling evaluation in real mobile device scenarios, testing temporal
parsing, implicit user preference, and contextual instruction following.

\subsection{2025: Dual Control, Reinforcement Learning, MCP, and International Scope}

\paragraph{\tautwobench{} \citep{barres2025tau2}.}
Formalised the dual-control setting as a Dec-POMDP. In the telecom domain (agent CRM tools + user
phone tools, 30 total), the best model (claude-3.7-sonnet) achieves only 49\% pass@1, dropping
below 25\% at pass@4.

\paragraph{ReTool \citep{feng2025retool}.}
Applied reinforcement learning to tool-use, training models to decide when tool calling is more
efficient than chain-of-thought reasoning alone. Demonstrates significant gains on multi-step
mathematical reasoning.

\paragraph{MCPVerse / MCPToolBench \citep{mcpverse2025}.}
Evaluation over real-world APIs via the Model Context Protocol (MCP, Anthropic 2024), the emerging
standard for tool discovery and context management. The largest real-API benchmark to date.

\paragraph{International Tool Calling (ITC) \citep{itc2025}.}
A large-scale multilingual benchmark covering 3,571 real APIs and 17,540 tool-calling tasks across
20 API categories and 40 countries. Fine-tuning on the full multilingual ITC dataset improves
cross-lingual generalisation substantially: tool selection recall improves by up to 42\% on
non-English tasks versus English-only training (22.6\% improvement). Among open-source models,
DeepSeek-V3 and Qwen2.5-Coder-32B perform strongly; Watt-tool-8B scores high on task accuracy but
poorly on language matching, revealing weaknesses in multilingual structural adherence.

\paragraph{MASSIVE-Agents \citep{massiveagents2025}.}
The first truly multilingual function-calling benchmark covering 51 languages including Hindi, built
by converting the MASSIVE NLU dataset into function-call format. GPT-4-class models achieve
$\sim$50--60\% in Hindi versus $\sim$80\% in English; smaller open-source models fall to
15--30\% in Hindi.

\paragraph{ComplexFuncBench \citep{zhong2025complexfuncbench}.}
Evaluates nested calls, constraint-aware argument generation, and context-window management under
long-context scenarios.

\paragraph{When2Call \citep{when2call2025}.}
Explicitly tests the model's decision among: (a) call a tool, (b) ask for more info, (c) say it
cannot help, or (d) answer directly. Fills the gap between \bfcl{}'s binary irrelevance test and
real production behaviour.

\paragraph{Pre-Act \citep{preact2025}.}
Structured multi-step planning before action generation. Plans with per-step reasoning outperform
ReAct in agentic settings.

% ══════════════════════════════════════════════════════════════════════════
\section{Benchmarks and Evaluation}
\label{sec:benchmarks}
% ══════════════════════════════════════════════════════════════════════════

\subsection{Benchmark Landscape}

Table~\ref{tab:benchmarks} provides a systematic comparison of major benchmarks across the key
design dimensions introduced in Section~\ref{sec:taxonomy}.

\begin{table}[h]
\centering
\caption{Systematic comparison of major tool-calling benchmarks (2023--2025).
\yes~= yes, \no~= no, $\sim$~= partial.}
\label{tab:benchmarks}
\renewcommand{\arraystretch}{1.2}
\small
\begin{tabular}{L{2.4cm} C{0.6cm} L{1.7cm} L{1.6cm} C{0.9cm} C{0.9cm} C{0.7cm} L{2.0cm}}
\toprule
\rowcolor{headerbg}
\textbf{Benchmark} & \textbf{Year} & \textbf{Scale} & \textbf{Eval Method}
  & \textbf{Multi-Turn} & \textbf{User Sim.} & \textbf{Policy} & \textbf{Multilingual} \\
\midrule
APIBench (Gorilla)   & 2023 & 1,645 APIs     & AST         & \no  & \no      & \no  & \no \\
\rowcolor{altrow}
API-Bank             & 2023 & 73 APIs        & Execution   & \yes & \no      & \no  & \no \\
ToolBench            & 2023 & 16K+ APIs      & LLM judge   & \yes & \no      & \no  & \no \\
\rowcolor{altrow}
\bfcl{} v1           & 2024 & 2,000 pairs    & AST+Exec    & \no  & \no      & \no  & \no \\
\bfcl{} v3           & 2024 & 5,551 pairs    & AST+State   & \yes & \no      & \no  & \no \\
\rowcolor{altrow}
\taubench{}          & 2024 & 165 tasks      & State (DB)  & \yes & LLM sim  & \yes & \no \\
ToolSandbox          & 2024/25 & Custom      & Milestone   & \yes & LLM sim  & $\sim$ & \no \\
\rowcolor{altrow}
ITC Dataset          & 2025 & 17,540 tasks   & AST+Exec    & \no  & \no      & \no  & \yes (40 countries) \\
\tautwobench{}       & 2025 & 114 tasks      & Assertions  & \yes & LLM sim  & \yes & \no \\
\rowcolor{altrow}
MASSIVE-Agents       & 2025 & 51 langs       & AST         & \no  & \no      & \no  & \yes (incl.\ Hindi) \\
Arabic \bfcl{}       & 2025 & \bfcl{}-scale  & AST         & \no  & \no      & \no  & Arabic \\
\rowcolor{altrow}
ACEBench             & 2025 & Multi-domain   & Rule+LLM    & \yes & Partial  & \no  & Chinese+EN \\
\bottomrule
\end{tabular}
\end{table}

\subsection{Evaluation Methods Compared}

\paragraph{AST Evaluation.}
Fast, scalable, and deterministic with no LLM-judge bias. Cannot catch semantic errors where the
call is syntactically valid but semantically wrong; requires exhaustive labelling of synonym values.

\paragraph{Executable Evaluation.}
Catches semantic errors; ground truth is the execution result. Requires real or simulated API
implementations. API stability is a persistent challenge, addressed by
StableToolBench~\citep{guo2024stabletoolbench}.

\paragraph{State-Based Evaluation.}
Checks outcomes rather than calls; path-agnostic (alternative valid action sequences pass). Requires
a full simulation environment and does not penalise inefficient paths.

\paragraph{LLM-as-Judge.}
Flexible and handles open-ended outputs, but exhibits GPT-family bias, verbosity bias, and is not
reproducible across model versions. ToolBench's ChatGPT-based ToolEval has been widely criticised on
these grounds.

\subsection{Key Empirical Findings Across Benchmarks}

\begin{enumerate}[leftmargin=1.4em, itemsep=3pt]
  \item \textbf{Simple single-turn is near-saturated.} Top models exceed 90\% on \bfcl{} simple
        Python. This category no longer discriminates frontier models.
  \item \textbf{Parallel and multiple function calls remain hard.} A 70--80 percentage-point gap
        exists between proprietary and open-source models on \bfcl{} parallel-multiple.
  \item \textbf{Multi-turn is the real frontier.} \bfcl{} composite multi-turn and \taubench{}
        both show 40--60~pp drops from simple single-turn performance.
  \item \textbf{Dual control is dramatically harder.} \tautwobench{} telecom drops 22--40~pp
        versus retail/airline for the same model family.
  \item \textbf{Consistency collapses at $k \geq 4$.} GPT-4o: 61\% pass@1 $\rightarrow$
        $<$25\% pass@8 on \taubench{} retail.
  \item \textbf{Irrelevance detection is a persistent failure.} Even GPT-4o generates hallucinated
        tool calls when no tool is appropriate ($\sim$15--25\% error rate).
  \item \textbf{Native FC mode outperforms prompt injection.} Models using native function-calling
        APIs consistently outperform text-injection approaches on the same base model.
  \item \textbf{Cross-lingual generalisation requires multilingual training.} ITC~\citep{itc2025}
        demonstrates that restricting training to English reduces non-English tool selection recall
        by up to 42\% compared to full multilingual training.
\end{enumerate}

% ══════════════════════════════════════════════════════════════════════════
\section{Training Methods for Tool Calling}
\label{sec:training}
% ══════════════════════════════════════════════════════════════════════════

\subsection{Supervised Fine-Tuning (SFT)}

SFT remains the dominant training approach. Three primary data sources are used:

\textbf{Distillation from closed models.} GPT-4 generates (query, tool\_call) pairs; open-source
models are fine-tuned on these. Used by ToolLLM, ToolAlpaca, and early Gorilla. Licensing
restrictions limit commercial deployment.

\textbf{Synthetic data with self-evolved APIs.} ToolACE's TSS generates tools from scratch; quality
is verified by rule checkers and model checkers. Avoids licensing issues entirely.

\textbf{Real API data.} ToolBench scraped RapidAPI. Quality issues ($\approx$50\% of queries and
75\% of trajectories had errors in the original release) were addressed by StableToolBench and
ToolBench-V.

Key SFT challenges include: (1)~\emph{catastrophic forgetting} --- fine-tuning for tool use
degrades general capability, mitigated by mixing general instruction data; (2)~\emph{data quality
vs.\ quantity} --- 10K high-quality examples outperform 100K noisy ones; and
(3)~\emph{complexity calibration} --- training data slightly above model capability improves
generalisation (the self-guided approach in ToolACE).

\subsection{Reinforcement Learning}

2025 saw a major shift from SFT to RL for tool calling.
ReTool~\citep{feng2025retool} trains models to decide when tool calling is more efficient than
chain-of-thought reasoning alone, achieving significant gains on multi-step mathematical reasoning.
ToolACE-R uses self-refinement based on internal confidence signals without external oracle feedback.
TRICE-style second-stage RL uses execution success as a reward signal.

The key advantage of RL is its ability to discover non-obvious tool-call strategies absent from
training demonstrations, particularly effective for multi-step and long-horizon tasks.

\subsection{Retrieval-Augmented Tool Selection}

When the tool pool spans thousands of APIs, embedding-based retrieval selects candidate APIs before
passing them to the LLM. ToolLLM's neural API retriever achieves 90\%+ retrieval precision on
ToolBench. Progressive retrieval (ProTIP) uses multi-stage contrastive training for context-aware
tool selection. Gorilla's RAT jointly fine-tunes retriever and generator so the generator produces
calls compatible with retrieved API documentation.

\emph{Challenge:} Long-context models with 128K+ context windows may reduce the need for retrieval,
but introduce needle-in-haystack accuracy degradation for large tool pools.

% ══════════════════════════════════════════════════════════════════════════
\section{Failure Modes and Error Taxonomies}
\label{sec:failures}
% ══════════════════════════════════════════════════════════════════════════

\subsection{Single-Turn Failures}

From \bfcl{} and related benchmarks, failures cluster into seven categories:
(1)~wrong function name; (2)~wrong argument name; (3)~wrong argument value;
(4)~wrong argument type (string vs.\ integer, list vs.\ scalar); (5)~hallucinated arguments
(parameters invented that do not exist in schema); (6)~missing required arguments; and
(7)~irrelevance hallucination (a call generated when no function is applicable).

\subsection{Multi-Turn Failures}

From \taubench{} trajectory analysis (GPT-4o on $\tau$-retail, $n=36$):
\begin{itemize}[leftmargin=1.2em, itemsep=2pt]
  \item \textbf{Wrong argument (33\%):} Correct API, wrong values --- often from failure to reason
        over complex nested database structures.
  \item \textbf{Wrong decision / policy violation (25\%):} Calls the wrong API by failing to apply
        domain policy (\eg calling \texttt{exchange} twice when policy allows only once).
  \item \textbf{Wrong info (22\%):} Omits required output information or provides incorrect totals.
  \item \textbf{Partial resolution (19\%):} Handles some but not all user requests.
\end{itemize}

\subsection{Multilingual Failures}

From \emph{Lost in Execution}~\citep{lostinexecution2025}, MASSIVE-Agents~\citep{massiveagents2025},
and ITC~\citep{itc2025}:
\begin{itemize}[leftmargin=1.2em, itemsep=2pt]
  \item \textbf{Cross-lingual schema mismatch:} When function descriptions are in English but
        queries are in non-English scripts, argument extraction error rates increase substantially.
  \item \textbf{Parameter hallucination:} Higher hallucination rates for non-English queries even
        when English function schemas are used.
  \item \textbf{Tokenisation inefficiency:} Indic languages exhibit high token fertility
        (3--6 subwords per word in standard tokenisers), reducing effective context for tool schemas.
  \item \textbf{Execution-level violations:} Output fails formatting requirements specific to
        non-English scripts (mixed-direction text, script-specific tokenisation).
  \item \textbf{Structural non-adherence:} Models like Watt-tool-8B achieve strong task-level
        performance but poor language matching (LM) and format matching (FM) in multilingual
        settings~\citep{itc2025}.
\end{itemize}

% ══════════════════════════════════════════════════════════════════════════
\section{Multilingual Tool Calling: The State of the Art}
\label{sec:multilingual}
% ══════════════════════════════════════════════════════════════════════════

\subsection{MASSIVE-Agents and the Multilingual Frontier}

MASSIVE-Agents~\citep{massiveagents2025} is the first truly multilingual function-calling benchmark
covering 51 languages including Hindi. Built on the MASSIVE NLU dataset, each intent becomes a
function name and each slot becomes a parameter; function descriptions are maintained in English
while queries appear in the target language. GPT-4-class models achieve approximately 50--60\% in
Hindi versus 80\% in English; smaller open-source models fall to 15--30\% in Hindi compared to
50--60\% in English.

MASSIVE-Agents is nonetheless limited to a slot-filling style derived from NLU data, not from real
production APIs. Multi-turn scenarios, parallel calling, and policy adherence are absent.

\subsection{International Tool Calling (ITC)}

The ITC dataset~\citep{itc2025} is the most geographically diverse tool-calling benchmark to date.
It covers 3,571 real APIs and 17,540 tasks across 20 API categories and 40 countries, constructed
through a four-stage pipeline: API collection, query generation, scoring/filtering, and QA pair
generation. Key findings:
\begin{itemize}[leftmargin=1.2em, itemsep=2pt]
  \item Among closed-source models, GPT-4o achieves 97.95\% language matching (LM); DeepSeek-R1
        achieves 100\% format matching (FM).
  \item Models trained on the full multilingual ITC dataset improve tool selection recall by 42\%
        on non-English tasks, versus 22.6\% for English-only training.
  \item Tool invocation F1 for Qwen2.5-Coder-7B increased 34.6\% with multilingual training,
        outperforming the English-only gain by 9.3\%.
\end{itemize}

\subsection{Language-Specific Benchmarks}

\textbf{Arabic \bfcl{} Extension (2025):} Translates \bfcl{} queries to Arabic and evaluates
open-source Arabic LLMs. All models show significant degradation versus English, attributable to
tokenisation inefficiency, insufficient Arabic tool-use training data, and suboptimal internal
representations~\citep{arabicbfcl2025}.

\textbf{FunctionChat-Bench (Korean):} Multi-turn Korean function-calling evaluation. Demonstrates
that language-specific benchmarks are necessary: Korean-specific models outperform generic
multilingual models on Korean function calling.

\textbf{ACEBench (Chinese + English):} Systematic robustness assessment across diverse scenarios
with rule-based and LLM-based evaluation.

\subsection{Structural Challenges for Non-English Tool Calling}

\begin{itemize}[leftmargin=1.2em, itemsep=2pt]
  \item \textbf{Token fertility:} Indic and Arabic scripts have 3--6 subwords per word in standard
        BPE tokenisers, effectively reducing the context window available for tool schemas by a
        comparable factor.
  \item \textbf{Morphological complexity:} Agglutinative morphology in Dravidian languages (Tamil,
        Telugu, Kannada, Malayalam) makes entity boundary detection harder.
  \item \textbf{Script diversity:} Nine distinct scripts across India's scheduled languages, plus
        code-switching (Hinglish, Tanglish), create unpredictable tokenisation patterns.
  \item \textbf{Training data scarcity:} No publicly available (query, function\_schema,
        ground\_truth\_call) triples exist in any Indic language, creating a circular dependency
        between model quality and data availability.
\end{itemize}

% ══════════════════════════════════════════════════════════════════════════
\section{The Indic Tool Calling Frontier}
\label{sec:indic}
% ══════════════════════════════════════════════════════════════════════════

\subsection{The Indic LLM Ecosystem}

\paragraph{AI4Bharat (IIT Madras).}
AI4Bharat has produced IndicBERT, IndicBART, IndicTrans2 (110+ translation directions across all
22 scheduled languages), Airavata~\citep{gala2024airavata} (Hindi instruction-tuned model),
and IndicLLMSuite~\citep{indicllmsuite2024}. Research focus has been on NLU, translation, and
generation --- no tool-calling capability has been published.

\paragraph{Sarvam AI.}
Sarvam has released OpenHathi (a 7B Hindi-English bilingual model based on LLaMA-2), Sarvam-1
\citep{sarvam12024} (a 2B parameter model supporting 10 Indic languages, trained on 4T tokens with
a custom tokeniser reducing token fertility), and Sarvam-M (a 24B multilingual reasoning model).
No tool-calling evaluation has been published.

\paragraph{Krutrim (Ola).}
Krutrim-1 (7B Mistral-based, 2T+ tokens) and Krutrim-2 (13 languages) represent the proprietary
Indic LLM track. The \textit{Kruti} assistant (2025) --- booking rides, paying bills, ordering
food across 13 Indian languages --- is the closest production deployment to Indic tool calling, but
no academic benchmark has been published.

\paragraph{Other Ecosystem Players.}
Tech Mahindra's Project Indus processes $\sim$100M queries/month. Two AI's SUTRA supports 50+
languages. Academic monolingual models include Tamil-LLaMA, Kannada LLaMA, and Odia-Llama.
IndicGenBench~\citep{indicgenbench2024} (ACL 2024) covers 29 Indic languages for generation tasks.
IndicMMLU-Pro~\citep{indicmmlupro2025} benchmarks nine Indic languages on academic knowledge.
IBM's MILU~\citep{milu2025} covers 11 Indic languages across 41 subjects from Indian competitive
examinations. The Indic-QA Benchmark~\citep{indicqa2025} (NAACL 2025) targets context-grounded
question answering in 11 Indic languages. \emph{None of these includes tool-calling tasks.}

\subsection{Real-World Indic Tool-Calling Deployments}

\begin{itemize}[leftmargin=1.2em, itemsep=3pt]
  \item \textbf{BharatGPT / CoRover.ai:} Deployed for IRCTC (Indian Railways), banking, and
        government services. Claims 1B+ users across Hindi, Tamil, Bengali, and other Indic
        languages for booking, enquiry, and transaction workflows. Proprietary; no academic benchmark.
  \item \textbf{Kruti (Krutrim):} Agentic assistant in 13 Indian languages --- the first
        commercially deployed Indic agent. No academic benchmark paper.
  \item \textbf{Sarvam + Aadhaar (UIDAI):} Voice-based multilingual AI for identity verification
        queries in regional languages.
  \item \textbf{India AI Mission:} Rs.\,10,371 crore investment with 27 AI labs in Tier~2/3
        cities. ``Citizen Connect 2047'' and ``AI4Pragati'' use cases require multilingual agentic
        AI at scale.
\end{itemize}

\subsection{Gap Analysis: What Exists vs.\ What Is Missing}

\textbf{What exists:}
MASSIVE-Agents Hindi component (slot-filling style, single-turn only);
\emph{Lost in Execution}~\citep{lostinexecution2025} (10 languages including Hindi, single-turn
only); ITC dataset (real APIs, 40 countries, but no explicit Indic language subsection or
India-specific API domains).

\textbf{What is missing:}
(1)~No Indic tool-calling benchmark covering real India-specific APIs (UPI, Aadhaar, IRCTC, mandi
prices). (2)~No multi-turn Indic benchmark. (3)~No code-switching (Hinglish/Tanglish) evaluation.
(4)~No voice-to-tool pipeline evaluation for low-literacy users. (5)~No systematic evaluation of
any Indic LLM on tool calling. (6)~No study of culturally-specific API design as function-calling
schemas.

\subsection{Proposed Benchmark: \textsc{IndicFuncBench}}

\begin{researchbox}{IndicFuncBench Specification}
We propose \textsc{IndicFuncBench}, the first function-calling benchmark designed for Indic
languages, covering:
\begin{itemize}[leftmargin=1.0em, itemsep=2pt, topsep=2pt]
  \item \textbf{Languages:} Hindi, Bengali, Tamil, Telugu, Kannada, Malayalam, Gujarati, Marathi,
        Odia, Punjabi (10 languages, 4 script families).
  \item \textbf{Domains:} (1)~Financial (UPI transfer, EMI calculator, bank balance);
        (2)~Government services (Aadhaar lookup, ration card, PAN verification);
        (3)~Commerce (railway PNR, bus booking, grocery delivery);
        (4)~Agriculture/Rural (mandi price query, crop disease detection, kisan credit).
  \item \textbf{Categories:} Simple $\cdot$ Multiple $\cdot$ Parallel $\cdot$ Irrelevance
        $\cdot$ Multi-Turn (mirroring \bfcl{} structure).
  \item \textbf{Evaluation:} Executable evaluation + AST; human annotation by native speakers;
        not machine-translated from English.
  \item \textbf{Key design choices:} Function schemas in English (matching production);
        code-switched queries (Hinglish, Tanglish) included explicitly;
        pass@$k$ with $k \geq 4$ for consistency measurement.
\end{itemize}
\end{researchbox}

% ══════════════════════════════════════════════════════════════════════════
\section{Open Problems and Future Directions}
\label{sec:openproblems}
% ══════════════════════════════════════════════════════════════════════════

\subsection{Evaluation Gaps}

\textbf{Safety-aware tool calling.} What happens when a model calls the wrong tool with financial
or medical consequences? Safety-critical evaluation (privacy leaks, wrong transactions, medical
contraindications) remains largely unexplored.

\textbf{Long-context tool calling.} As context windows grow to 1M tokens, how does multi-step tool
sequence performance change? ComplexFuncBench~\citep{zhong2025complexfuncbench} begins to address
this.

\textbf{Dynamic tool creation.} When no existing tool can answer a query, can a model generate a
new tool specification? UltraTool explores this but it remains underexplored.

\textbf{MCP ecosystem evaluation.} With Model Context Protocol adoption, evaluation must move from
simulated to real-world API calls. MCPVerse and MCPBench are early steps.

\subsection{Training Gaps}

\textbf{RL from real execution.} Most RL work uses simulated environments. Training on real API
outcomes with sandboxing remains an open problem.

\textbf{Sample efficiency.} Current fine-tuning requires thousands of examples. Few-shot
in-context tool learning for new domains without fine-tuning is underdeveloped.

\textbf{Compositional generalisation.} Models trained on individual tools often fail when two tools
must be composed in a novel way.

\textbf{Forgetting mitigation at scale.} Fine-tuning for tool calling degrades general capability.
Principled solutions beyond data mixing are needed.

\subsection{Indic-Specific Research Priorities}

\begin{enumerate}[leftmargin=1.4em, itemsep=3pt]
  \item \textbf{IndicFuncBench:} Create the first public Indic tool-calling benchmark
        (Section~\ref{sec:indic}).
  \item \textbf{Tokeniser improvement:} Indic-script-aware BPE reducing token fertility.
        Sarvam-1's custom tokeniser is a step forward.
  \item \textbf{Tool-calling SFT data:} Create (query, schema, call) triples for 10 Indic
        languages, starting with translation + human verification.
  \item \textbf{Code-switching robustness:} Explicitly include Hinglish/Tanglish training
        examples; real users code-switch extensively.
  \item \textbf{Voice-to-tool pipelines:} For low-literacy users, the practical pipeline is
        ASR $\rightarrow$ language ID $\rightarrow$ normalisation $\rightarrow$ tool call.
        IndicWav2Vec + Sarvam-1 + function calling is an unexplored combination.
\end{enumerate}

% ══════════════════════════════════════════════════════════════════════════
\section{Related Surveys}
\label{sec:related}
% ══════════════════════════════════════════════════════════════════════════

\textbf{Tool-calling surveys.}
\citet{wangsurvey2024} and the Springer 2025 comprehensive survey~\citep{springersurvey2025} cover
SFT, RL, and retrieval approaches systematically, but are entirely English-centric.

\textbf{Agent evaluation surveys.}
\citet{agentsurvey2025} provide a broad agent evaluation survey including tool use, but with no
multilingual analysis.

\textbf{Indic LLM surveys.}
\citet{indicanalysis2025} review 14 papers evaluating LLMs on Indic languages, covering no tool
calling. IndicMMLU-Pro~\citep{indicmmlupro2025} and the Indic-QA Benchmark~\citep{indicqa2025}
represent the frontier of Indic NLU evaluation but neither touches tool calling.

This survey is the first to bridge all three, providing a complete map of the global tool-calling
literature and a concrete roadmap for Indic tool calling.

% ══════════════════════════════════════════════════════════════════════════
\section{Conclusion}
\label{sec:conclusion}
% ══════════════════════════════════════════════════════════════════════════

Tool calling has emerged as the foundational capability for agentic AI, evolving rapidly from
Toolformer's self-supervised five-tool system in 2023 to today's ecosystem of dual-control agents,
reinforcement-learned decision making, real-API evaluation over the Model Context Protocol, and
internationally diverse benchmarks like ITC spanning 40 countries. The frontier has shifted from
simple single-turn function selection --- now near-saturated at 90\%+ for top models --- to
multi-turn consistency, policy adherence, and cross-lingual robustness.

Yet 1.4 billion speakers of India's 22 scheduled languages remain almost entirely unserved by
existing tool-calling research. The Indic LLM ecosystem has produced capable multilingual foundation
models (Sarvam-1, Krutrim-2, Airavata), and production deployments (Kruti, BharatGPT) demonstrate
clear demand. What is missing is the evaluation infrastructure: no tool-calling benchmark covers
India-specific APIs in Indic languages, no systematic comparison of Indic LLMs on tool calling
exists, and no work addresses the full voice-to-tool pipeline critical for rural and low-literacy
users.

We present \textsc{IndicFuncBench} as a concrete specification for addressing this gap, designed to
cover 10 Indic languages, four India-specific domain clusters, and the full \bfcl{} evaluation
category spectrum. We hope this survey serves as both a reference for researchers entering the field
and a catalyst for the Indic tool-calling research community.

% ══════════════════════════════════════════════════════════════════════════
% REFERENCES
% ══════════════════════════════════════════════════════════════════════════
\bibliographystyle{plainnat}
\begin{thebibliography}{99}

\bibitem[Barres et al.(2025)]{barres2025tau2}
Barres, V. et al.\ (2025).
\newblock $\tau^2$-bench: Evaluating Conversational Agents in a Dual-Control Environment.
\newblock \textit{arXiv:2506.07982}.

\bibitem[Feng et al.(2025)]{feng2025retool}
Feng, Y. et al.\ (2025).
\newblock ReTool: Reinforcement Learning for Strategic Tool Use in LLMs.
\newblock \textit{arXiv:2504.11536}.

\bibitem[Gala et al.(2024)]{gala2024airavata}
Gala, J. et al.\ (2024).
\newblock Airavata: Introducing Hindi Instruction-tuned LLM.
\newblock \textit{arXiv:2401.15006}.

\bibitem[Guo et al.(2024)]{guo2024stabletoolbench}
Guo, L. et al.\ (2024).
\newblock StableToolBench: Towards Stable Large-Scale Benchmarking on Tool Learning of Large Language Models.
\newblock \textit{arXiv:2403.07714}.

\bibitem[Huang et al.(2023)]{huang2023metatool}
Huang, Y. et al.\ (2023).
\newblock MetaTool Benchmark for Large Language Models: Deciding Whether to Use Tools and Which to Use.
\newblock \textit{arXiv:2310.03128}.

\bibitem[ITC(2025)]{itc2025}
Anonymous (2025).
\newblock Enhancing Tool Calling in LLMs with the International Tool Calling Dataset.
\newblock \textit{ICLR 2026 Submission, arXiv:2603.05515}.

\bibitem[Kunchukuttan et al.(2024)]{indicllmsuite2024}
AI4Bharat (2024).
\newblock IndicLLMSuite: A Blueprint for Creating Pre-training and Fine-Tuning Datasets for Indian Languages.
\newblock \textit{arXiv:2403.06350}.

\bibitem[Li et al.(2023)]{li2023apibank}
Li, M. et al.\ (2023).
\newblock API-Bank: A Comprehensive Benchmark for Tool-Augmented LLMs.
\newblock In \textit{EMNLP 2023}.

\bibitem[Liu et al.(2024)]{liu2024xlam}
Liu, W. et al.\ (2024).
\newblock xLAM: A Family of Large Action Models to Empower AI Agent Systems.
\newblock \textit{arXiv:2409.03215}.

\bibitem[Liu et al.(2025)]{liu2025toolace}
Liu, W. et al.\ (2025).
\newblock ToolACE: Winning the Points of LLM Function Calling.
\newblock In \textit{ICLR 2025}.

\bibitem[Lu et al.(2025)]{lu2025toolsandbox}
Lu, J. et al.\ (2025).
\newblock ToolSandbox: A Stateful, Conversational, Interactive Evaluation Benchmark for LLM Tool Use Capabilities.
\newblock In \textit{NAACL 2025}.

\bibitem[MASSIVE-Agents(2025)]{massiveagents2025}
Anonymous (2025).
\newblock MASSIVE-Agents: A Benchmark for Multilingual Function Calling.
\newblock In \textit{ACL 2025 Findings}.

\bibitem[MCPVerse(2025)]{mcpverse2025}
Anonymous (2025).
\newblock MCPVerse: Evaluation of LLMs over Real-World APIs via Model Context Protocol.
\newblock \textit{arXiv:2025}.

\bibitem[LostInExecution(2025)]{lostinexecution2025}
Anonymous (2025).
\newblock Lost in Execution: On the Multilingual Robustness of Tool Calling in LLMs.
\newblock \textit{arXiv:2601.05366}.

\bibitem[ArabicBFCL(2025)]{arabicbfcl2025}
Anonymous (2025).
\newblock Arabic Prompts with English Tools: A Function Calling Benchmark.
\newblock \textit{arXiv:2601.05101}.

\bibitem[Parisi et al.(2022)]{parisi2022talm}
Parisi, A., Zhao, Y., and Fiedel, N. (2022).
\newblock TALM: Tool Augmented Language Models.
\newblock \textit{arXiv:2205.12255}.

\bibitem[Patil et al.(2023)]{patil2023gorilla}
Patil, S.~G. et al.\ (2023).
\newblock Gorilla: Large Language Model Connected with Massive APIs.
\newblock In \textit{NeurIPS 2024}.

\bibitem[Patil et al.(2025)]{patil2025bfcl}
Patil, S.~G. et al.\ (2025).
\newblock The Berkeley Function Calling Leaderboard.
\newblock In \textit{ICML 2025, PMLR 267}.

\bibitem[PreAct(2025)]{preact2025}
Anonymous (2025).
\newblock Pre-Act: Multi-Step Planning Before Action in Agentic LLMs.
\newblock \textit{arXiv:2025}.

\bibitem[Qin et al.(2024)]{qin2024toolllm}
Qin, Y. et al.\ (2024).
\newblock ToolLLM: Facilitating Large Language Models to Master 16000+ Real-world APIs.
\newblock In \textit{ICLR 2024}.

\bibitem[Sarvam(2024)]{sarvam12024}
Sarvam AI (2024).
\newblock Sarvam-1: A Language Model for Indian Languages.
\newblock Technical Report.

\bibitem[Schick et al.(2023)]{schick2023toolformer}
Schick, T. et al.\ (2023).
\newblock Toolformer: Language Models Can Teach Themselves to Use Tools.
\newblock In \textit{NeurIPS 2023}.

\bibitem[Shen et al.(2023)]{shen2023hugginggpt}
Shen, Y. et al.\ (2023).
\newblock HuggingGPT: Solving AI Tasks with ChatGPT and its Friends in Hugging Face.
\newblock In \textit{NeurIPS 2023}.

\bibitem[Singh et al.(2024)]{indicgenbench2024}
Singh, H. et al.\ (2024).
\newblock IndicGenBench: A Multilingual Benchmark to Evaluate Generation Capabilities of LLMs on Indic Languages.
\newblock In \textit{ACL 2024}.

\bibitem[Singh et al.(2025)]{indicqa2025}
Singh, A.~K. et al.\ (2025).
\newblock INDIC QA BENCHMARK: A Multilingual Benchmark to Evaluate Question Answering Capability of LLMs for Indic Languages.
\newblock In \textit{NAACL 2025 Findings}.

\bibitem[Srinivasan et al.(2023)]{srinivasan2023nexusraven}
Srinivasan, V. et al.\ (2023).
\newblock NexusRaven: A Commercially-Permissive Language Model for Function Calling.
\newblock In \textit{NeurIPS 2023 Workshop}.

\bibitem[Tang et al.(2023)]{tang2023toolalpaca}
Tang, Q. et al.\ (2023).
\newblock ToolAlpaca: Generalized Tool Learning for Language Models with 3000 Simulated Cases.
\newblock \textit{arXiv:2306.05301}.

\bibitem[Thoppilan et al.(2022)]{thoppilan2022lamda}
Thoppilan, R. et al.\ (2022).
\newblock LaMDA: Language Models for Dialog Applications.
\newblock \textit{arXiv:2201.08239}.

\bibitem[IndicAnalysis(2025)]{indicanalysis2025}
Anonymous (2025).
\newblock Analysis of Indic Language Capabilities in LLMs.
\newblock \textit{arXiv:2501.13912}.

\bibitem[IndicMMLUPro(2025)]{indicmmlupro2025}
Anonymous (2025).
\newblock IndicMMLU-Pro: Benchmarking Indic Large Language Models on Multi-Task Language Understanding.
\newblock \textit{arXiv:2501.15747}.

\bibitem[MILU(2025)]{milu2025}
Verma, S. et al.\ (2025).
\newblock MILU: A Multi-Task Indic Language Understanding Benchmark.
\newblock In \textit{NAACL 2025}.

\bibitem[Wang et al.(2024)]{wangsurvey2024}
Wang, L. et al.\ (2024).
\newblock Tool Learning with Large Language Models: A Survey.
\newblock \textit{arXiv:2405.17935}.

\bibitem[Wang et al.(2024)]{wang2024hammerbench}
Wang, Y. et al.\ (2024).
\newblock HammerBench: Fine-Grained Function-Calling Evaluation in Real Mobile Device Scenarios.
\newblock \textit{arXiv:2410.04587}.

\bibitem[When2Call(2025)]{when2call2025}
Anonymous (2025).
\newblock When2Call: A Benchmark for Tool-Use Decision Making in LLMs.
\newblock \textit{arXiv:2025}.

\bibitem[Springer(2025)]{springersurvey2025}
Anonymous (2025).
\newblock Tool Learning with Language Models: A Comprehensive Survey of Methods, Pipelines, and Benchmarks.
\newblock \textit{Springer 2025}.

\bibitem[AgentSurvey(2025)]{agentsurvey2025}
Anonymous (2025).
\newblock Evaluation and Benchmarking of LLM Agents: A Survey.
\newblock \textit{arXiv:2507.21504}.

\bibitem[Yao et al.(2022a)]{yao2022webshop}
Yao, S. et al.\ (2022).
\newblock WebShop: Towards Scalable Real-World Web Interaction with Grounded Language Agents.
\newblock In \textit{NeurIPS 2022}.

\bibitem[Yao et al.(2023)]{yao2023react}
Yao, S. et al.\ (2023).
\newblock ReAct: Synergizing Reasoning and Acting in Language Models.
\newblock In \textit{ICLR 2023}.

\bibitem[Yao et al.(2024)]{yao2024taubench}
Yao, S. et al.\ (2024).
\newblock $\tau$-bench: A Benchmark for Tool-Agent-User Interaction in Real-World Domains.
\newblock \textit{arXiv:2406.12045}.      

\bibitem[Zhong et al.(2025)]{zhong2025complexfuncbench}
Zhong, W. et al.\ (2025).
\newblock ComplexFuncBench: Exploring Multi-Step and Constrained Function Calling under Long-Context Scenario.
\newblock \textit{arXiv:2501.10132}.

\end{thebibliography}

\end{document}